% Complete Section 7.4, original PDF 96-98 / printed 86-88.
% block: finite-element-differential-forms-heading page:96 type:section
\subsection{\BT{Finite element differential forms}{有限元微分形式}}\label{sec:finite-element-differential-forms}

% block: complete-form-degrees-of-freedom page:96 type:ordinary-paragraph
\BT{We are now ready to define our finite element spaces. As shape functions we use one of the complete polynomial spaces $\mathcal P_r\Lambda^k(T)$ or one of the trimmed polynomial spaces $\mathcal P_r^-\Lambda^k(T)$. An important observation is that the two spaces are related and are best understood if treated together. Indeed, the DOFs for the complete space involves the trimmed spaces on the faces and vice versa. Specifically, the DOFs for $\omega\in\mathcal P_r\Lambda^k(T)$ are}{现在可以定义我们的有限元空间. 形函数取自某个完全多项式空间 $\mathcal P_r\Lambda^k(T)$, 或某个修剪多项式空间 $\mathcal P_r^-\Lambda^k(T)$. 一个重要观察是, 两类空间彼此相关, 放在一起处理最易理解. 事实上, 完全空间的自由度涉及面上的修剪空间, 反之亦然. 具体地, $\omega\in\mathcal P_r\Lambda^k(T)$ 的自由度为}
\begin{equation}\label{eq:complete-form-degrees-of-freedom}
\omega\mapsto\int_f(\operatorname{tr}_f\omega)\wedge\mu,\qquad\mu\in\mathcal P_{r+k-d}^-\Lambda^{d-k}(f),\quad f\in\Delta_d(T),\ d\geq k,
\end{equation}

% block: trimmed-form-degrees-of-freedom page:97 type:ordinary-paragraph
\BT{while those for $\omega\in\mathcal P_r^-\Lambda^k(T)$ are}{而 $\omega\in\mathcal P_r^-\Lambda^k(T)$ 的自由度为}
\begin{equation}\label{eq:trimmed-form-degrees-of-freedom}
\omega\mapsto\int_f(\operatorname{tr}_f\omega)\wedge\mu,\qquad\mu\in\mathcal P_{r+k-d-1}\Lambda^{d-k}(f),\quad f\in\Delta_d(T),\ d\geq k.
\end{equation}
\BT{Note that each DOF is associated to a specific face $f$, of some dimension between $0$ and $n$, of the simplex $T$ ($\Delta_d(T)$ denotes the set of all faces of $T$ of dimension $d$). The assembled finite element space is then determined by the requirement that the DOFs are single-valued.}{注意, 每个自由度都与单纯形 $T$ 的某个特定面 $f$ 关联, 其维数介于 $0$ 与 $n$ 之间 ($\Delta_d(T)$ 表示 $T$ 的所有 $d$ 维面的集合). 然后要求自由度单值, 就确定了组装后的有限元空间.}

% block: form-unisolvence-goal page:97 type:ordinary-paragraph
\BT{In order that the DOF specifications \eqref{eq:complete-form-degrees-of-freedom} and \eqref{eq:trimmed-form-degrees-of-freedom} well define finite element spaces, we must verify \emph{unisolvence} for each of them. This requires showing, first, that the total number of DOFs we have associated to the shape function space is no more than its dimension and, second, that if all the DOFs vanish for some $\omega$ in the shape function space, then $\omega$ itself vanishes. This can be accomplished for both choices of shape function spaces. Here we verify unisolvence for the trimmed space $\mathcal P_r^-\Lambda^k(T)$, for any polynomial degree $r\geq1$, any dimension $n\geq1$, and any form degree $0\leq k\leq n$. The proof, which is based on the properties of the Koszul differential, is taken directly from \BookCite{6}. A very similar proof applies in the case of the complete polynomial spaces.}{为使自由度规定 \eqref{eq:complete-form-degrees-of-freedom} 和 \eqref{eq:trimmed-form-degrees-of-freedom} 良好地定义有限元空间, 必须分别验证其\emph{唯一可解性}. 这需要证明: 首先, 关联于形函数空间的自由度总数不超过空间维数; 其次, 若形函数空间中某个 $\omega$ 的所有自由度都为零, 则 $\omega$ 本身为零. 对两种形函数空间都可以做到这一点. 这里对任意多项式次数 $r\geq1$, 任意维数 $n\geq1$ 和任意形式次数 $0\leq k\leq n$, 验证修剪空间 $\mathcal P_r^-\Lambda^k(T)$ 的唯一可解性. 证明基于 Koszul 微分的性质, 直接取自 \BookCite{6}. 完全多项式空间的情形有十分相似的证明.}

% block: form-degree-of-freedom-count page:97 type:ordinary-paragraph
\BT{The number of DOFs defined in \eqref{eq:trimmed-form-degrees-of-freedom} is}{\eqref{eq:trimmed-form-degrees-of-freedom} 定义的自由度数为}
\[
\begin{aligned}
&\sum_{d\geq k}\#\Delta_d(T)\dim[\mathcal P_{r+k-d-1}\Lambda^{d-k}(\mathbb R^d)]\\
&\qquad=\sum_{d\geq k}\binom{n+1}{d+1}\binom{r+k-1}d\binom d{d-k}\\
&\qquad=\sum_{j\geq0}\binom{n+1}{j+k+1}\binom{r+k-1}{j+k}\binom{j+k}j.
\end{aligned}
\]
\BT{Simplifying with the binomial identity}{利用二项式恒等式化简,}
\[
\binom ab\binom bc=\binom ac\binom{a-c}{a-b}
\]
\BT{(which is immediate from the definition of the binomial coefficients), and Vandermonde's identity}{(它直接来自二项式系数的定义), 以及 Vandermonde 恒等式}
\[
\sum_{j\geq0}\binom a{b+j}\binom cj=\binom{a+c}{a-b},
\]
\BT{the right-hand side becomes}{右端变为}
\[
\binom{r+n}{r+k}\binom{r+k-1}k=\dim\mathcal P_r^-\Lambda^k(T),
\]
\BT{as required.}{如所要求.}

% block: form-unisolvence-induction pages:97,98 type:ordinary-paragraph
\BT{It remains to show that if all the DOFs vanish for some $\omega\in\mathcal P_r^-\Lambda^k(T)$, then $\omega$ vanishes. The proof proceeds by induction on the dimension $n$ (the case $n=1$ being easy). For any facet (face of codimension $1$) $f$ of $T$, $\operatorname{tr}_f\omega$ belongs to $\mathcal P_r^-\Lambda^k(f)$, and the DOFs for $\operatorname{tr}_f\omega$ vanish, since they are among the DOFs for $\omega$. Therefore we can invoke the inductive hypothesis to conclude that $\operatorname{tr}_f\omega$ vanishes for all facets $f$, i.e., that $\omega\in\mathring{\mathcal P}_r^-\Lambda^k(T)$, the subspace of elements of $\mathcal P_r^-\Lambda^k(T)$ whose traces vanish on the entire boundary. Since the trace operator commutes with the exterior derivative (as does any pullback), we have that $\nu:=d\omega\in\mathring{\mathcal P}_{r-1}\Lambda^{k+1}(T)$. Moreover,}{还需证明, 若某个 $\omega\in\mathcal P_r^-\Lambda^k(T)$ 的所有自由度都为零, 则 $\omega$ 为零. 对维数 $n$ 作归纳证明 (情形 $n=1$ 容易). 对 $T$ 的任意超面 (余维为 $1$ 的面) $f$, $\operatorname{tr}_f\omega$ 属于 $\mathcal P_r^-\Lambda^k(f)$, 且 $\operatorname{tr}_f\omega$ 的自由度为零, 因为它们是 $\omega$ 的自由度的一部分. 因此可用归纳假设得出, 对所有超面 $f$, $\operatorname{tr}_f\omega$ 都为零, 即 $\omega\in\mathring{\mathcal P}_r^-\Lambda^k(T)$, 这是 $\mathcal P_r^-\Lambda^k(T)$ 中在整个边界上迹为零的元素构成的子空间. 由于迹算子与外微分交换 (任意拉回都如此), 有 $\nu:=d\omega\in\mathring{\mathcal P}_{r-1}\Lambda^{k+1}(T)$. 此外,}
\[
\int_T\nu\wedge p=\pm\int_T\omega\wedge dp=0,\qquad p\in\mathcal P_{r+k-n}\Lambda^{n-k-1}(T),
\]
\BT{where the first equality comes from integration by parts and the second from the vanishing DOFs of $\omega$. From this one can show, by a concrete construction based on barycentric coordinates [\BookNumber{6}, Lemma 3.4], that $\nu=0$. Since $\omega\in\mathcal P_r^-\Lambda^k$ and $d\omega=0$, it follows that $\omega\in\mathcal P_{r-1}\Lambda^k$, which, together with the vanishing of its traces, means $\omega\in\mathring{\mathcal P}_{r-1}\Lambda^k$. A second application of the construction of [\BookNumber{6}, Lemma 3.4] shows that $\omega$ vanishes and completes the unisolvence proof.}{其中第一个等式来自分部积分, 第二个来自 $\omega$ 的自由度为零. 由此可以通过基于重心坐标的具体构造 [\BookNumber{6}, Lemma 3.4], 证明 $\nu=0$. 由于 $\omega\in\mathcal P_r^-\Lambda^k$ 且 $d\omega=0$, 可得 $\omega\in\mathcal P_{r-1}\Lambda^k$, 再结合其迹为零, 意味着 $\omega\in\mathring{\mathcal P}_{r-1}\Lambda^k$. 再次应用 [\BookNumber{6}, Lemma 3.4] 的构造, 表明 $\omega$ 为零, 从而完成唯一可解性的证明.}

% block: assembled-form-families page:98 type:ordinary-paragraph
\BT{Thus we have defined a family of spaces of polynomial differential $k$-forms on each element $T$ for use as shape functions, namely, the spaces $\mathcal P_r\Lambda^k(T)$ and $\mathcal P_r^-\Lambda^k(T)$, defined for each $r\geq1$. Moreover, for each of these spaces we have specified a unisolvent set of DOFs. Together with the triangulation $\mathcal T_h$, each such choice then determines a corresponding finite element space $\mathcal P_r\Lambda^k(\mathcal T_h)$ or $\mathcal P_r^-\Lambda^k(\mathcal T_h)$ of piecewise polynomial differential forms.}{由此在每个单元 $T$ 上定义了一族用作形函数的多项式微分 $k$-形式空间, 即对每个 $r\geq1$ 定义的空间 $\mathcal P_r\Lambda^k(T)$ 和 $\mathcal P_r^-\Lambda^k(T)$. 此外, 对每个空间都指定了一组具有唯一可解性的自由度. 结合剖分 $\mathcal T_h$, 每一种这样的选择都确定相应的分片多项式微分形式有限元空间 $\mathcal P_r\Lambda^k(\mathcal T_h)$ 或 $\mathcal P_r^-\Lambda^k(\mathcal T_h)$.}

% block: assembled-form-continuity-projections page:98 type:ordinary-paragraph
\BT{The DOFs defined in \eqref{eq:complete-form-degrees-of-freedom} and \eqref{eq:trimmed-form-degrees-of-freedom} have many nice properties, which are reflected in properties of the assembled finite element spaces. In particular,}{\eqref{eq:complete-form-degrees-of-freedom} 和 \eqref{eq:trimmed-form-degrees-of-freedom} 定义的自由度具有许多良好性质, 它们反映在组装后有限元空间的性质中. 特别地,}
\begin{itemize}
\item\BT{the DOFs enforce exactly the continuity needed for the assembled finite element function to belong to $H\Lambda^k$,}{这些自由度恰好施加了组装后的有限元函数属于 $H\Lambda^k$ 所需的连续性,}
\[
\mathcal P_r\Lambda^k(\mathcal T_h)=\{\,\omega\in H\Lambda^k(\Omega)\mid\omega|_T\in\mathcal P_r\Lambda^k(T)\,\};
\]
\item\BT{the canonical projection operators $\Pi_{r,h}^k:C\Lambda^k(\bar\Omega)\to\mathcal P_r\Lambda^k(\mathcal T_h)$ determined by the DOFs form a cochain map and similarly for the $\mathcal P^-$ spaces with their canonical projections $\Pi_{r,h}^{-,k}$.}{自由度确定的典范投影算子 $\Pi_{r,h}^k:C\Lambda^k(\bar\Omega)\to\mathcal P_r\Lambda^k(\mathcal T_h)$ 构成上链映射; 对 $\mathcal P^-$ 空间及其典范投影 $\Pi_{r,h}^{-,k}$ 也有类似结论.}
\end{itemize}

% block: canonical-projection-commutativity page:98 type:ordinary-paragraph
\BT{The commutativity of the canonical projections with the exterior derivative can be verified using the Stokes theorem. It is a crucial property: the main step toward the construction of bounded cochain projections.}{典范投影与外微分的交换性, 可用 Stokes 定理验证. 这是一个关键性质: 它是构造有界上链投影的主要步骤.}
